% Generated draft: requires mathematical and source review.
\documentclass[11pt]{article}
\usepackage[margin=25mm]{geometry}
\usepackage{fontspec}
\setmainfont{Latin Modern Roman}
\usepackage{amsmath,amssymb,mathtools,unicode-math}
\setmathfont{Latin Modern Math}
\usepackage{enumitem,needspace,xurl,hyperref}
\setlength{\parindent}{0pt}
\setlength{\parskip}{5pt}
\setlength{\emergencystretch}{3em}
\newcommand{\N}{\mathbb{N}}
\newcommand{\Z}{\mathbb{Z}}
\newcommand{\Q}{\mathbb{Q}}
\newcommand{\R}{\mathbb{R}}
\newcommand{\C}{\mathbb{C}}
\newcommand{\F}{\mathbb{F}}
\newcommand{\A}{\mathbb{A}}
\newcommand{\PP}{\mathbb{P}}
\newcommand{\E}{\mathbb{E}}
\newcommand{\eps}{\varepsilon}
\newcommand{\dd}{\,\mathrm{d}}
\begin{document}
\section*{UnsolvedMath ID: 20002675}
\textbf{Source problem code:} catalogue entry
\subsection{Definitions and mathematical statement}
\textbf{Question.} Find examples or conditions for free strong unimodality: a unimodal probability law \(\mu\) should remain unimodal after free additive convolution \(\mu\boxplus\nu\) with every unimodal law \(\nu\). Free convolution is the law of the sum of freely independent self-adjoint variables with the indicated laws. The source glosses unimodal as ``smooth density with a unique maximum'', giving Gaussian and semicircle examples. That wording does not fully fix the intended class: a unique global maximum alone permits smaller local peaks, while the semicircle density is not globally smooth across its support endpoints.

For the aid's conservative formulation, let \(\mathcal U\) be the standard unimodal probability laws, including the usual possible atom at the mode, with density nondecreasing to a mode and nonincreasing after it. Let \(\mathcal U_{\mathrm{reg}}\subseteq\mathcal U\) consist of laws with a globally \(C^\infty\) density and unique maximum. Standard free strong unimodality asks \(\mu\in\mathcal U\) and \(\mu\boxplus\nu\in\mathcal U\) for every \(\nu\in\mathcal U\). A smooth unique-mode version asks the analogous closure for \(\mu,\nu\in\mathcal U_{\mathrm{reg}}\). Failure even to belong to \(\mathcal U\) is enough to refute this stronger regular-output version.

\textbf{Reported finite-variance and smooth-witness obstructions.} The aid attributes to Hasebe and Ueda the theorem that no non-Dirac finite-variance probability measure is freely strongly unimodal in the standard sense. It additionally reports that a semicircle obstruction can be witnessed by some \(\nu\in\mathcal U_{\mathrm{reg}}\): a semicircle law \(\sigma\) satisfies
\[
 \sigma\boxplus\nu\notin\mathcal U.
\]
The actual witness density and full primary theorem references are not supplied. The smooth witness matters because a counterexample using only singular or nonsmooth test laws would not automatically settle a source reading restricted to smooth tests.

\textbf{Reported semicircular-attractor extension.} Suppose \(\mu\in\mathcal U_{\mathrm{reg}}\), and that nondegenerate affine normalizations of its free sums converge weakly to a semicircle:
\[
 D_{a_n}\bigl(\mu^{\boxplus n}\bigr)*\delta_{b_n}
 \ \Longrightarrow\ \sigma,\qquad a_n>0.
\]
Here \(D_a\) is dilation by \(a\), and the ordinary convolution with a point mass denotes translation. The aid reports that such \(\mu\) cannot satisfy the regular strong-unimodality condition. The conclusion uses the smooth semicircle witness and weak-limit/affine stability of standard unimodality; it is not a characterization of every infinite-variance law.

Its explicit example is
\[
 \mathrm d\mu(x)=Z^{-1}
 \frac{\mathrm dx}{(e+x^2)^{3/2}\log(e+x^2)},\qquad
 Z=\int_\R\frac{\mathrm dx}{(e+x^2)^{3/2}\log(e+x^2)}.
\]
This is a smooth symmetric probability density, uniquely maximized at zero and decreasing on the positive half-line. The aid reports finite first absolute moment but infinite variance. Its truncated second moment is unbounded and slowly varying; by Pata's criterion, the law belongs to the semicircular free domain of attraction. It therefore fails regular strong unimodality despite its infinite variance. The exact normalization sequence and full statement of Pata's criterion are not supplied.

The report says existence of any non-Dirac freely strongly unimodal law remains open in the literature checked. No date or complete bibliography accompanies that assessment. Dirac masses are trivial translation examples in the standard class but do not have smooth densities. The nonstandard unique-global-maximum-only interpretation is explicitly not separately settled by the supplied results. All reported obstructions require independent verification.
\subsection{Short English statement}
The source's unimodality convention needs clarification. The aid reports a standard finite-variance obstruction and strengthens the semicircle witness to a smooth uniquely peaked test law. A semicircular-attractor argument then excludes one explicit smooth infinite-variance example, without classifying all laws or settling a unique-global-maximum-only interpretation.
\subsection{Sources}
\begin{itemize}
\item Local UnsolvedMath registry, permanent ID \(20002675\), catalogue entry, \texttt{aim-probability-notes.json}, \texttt{notes[116]}, \emph{Free Analysis}, question 15, source dated 24 August 2006: \url{https://aimath.org/WWN/freeanalysis/freeanalysis.pdf}.
\item Supplied research attributions: Hasebe and Ueda on failure of non-Dirac finite-variance free strong unimodality; Pata's semicircular domain-of-attraction criterion. No titles, dates or full theorem hypotheses are supplied.
\item Supplied research summary on the smooth witness, semicircular-attractor extension and explicit infinite-variance density. Attempt 1: partial\_result, status at run ``open'', rigor ``proved'', low novelty confidence; registry status partially\_solved. The machine-generated aid is not peer reviewed and requires independent verification before citation; no PDF visit or updated literature search is claimed.
\end{itemize}
\end{document}
